\textbf{Tasks.} Five 6x9 mazes (start bottom, goal top right, reward 1 on reaching the goal, then reset to start; every episode shares one continuing budget of steps). \emph{Static}: a wall on row 2 with a gap at the right, $T=\rhval{const/horizon_steps}$. \emph{Blocking}: the gap is at the right until step \rhval{const/blocking_change_step}, then closes and a gap opens at the left, $T=\rhval{const/horizon_steps}$. \emph{Shortcut}: a longer path exists, and at step \rhval{const/shortcut_change_step} an additional shorter gap opens, $T=\rhval{const/shortcut_horizon_steps}$. \emph{Stochastic}: static layout, with probability \rhval{const/slip_prob:1} the chosen action is replaced by a uniformly random one. \emph{Stoch.\ blocking}: slip \rhval{const/slip_prob:1} and the blocking change together.

\textbf{Metrics.} \emph{cum} is the number of goal reaches over the run, \emph{early} over the first \rhval{const/reward_window_steps} steps, \emph{post} after the change (second half of the run for tasks without change) and \emph{late} over the last \rhval{const/reward_window_steps} steps. Higher is better for all.

\textbf{Hypotheses.} H1: on the static maze, planning raises reward over Q-learning (refuted if Dyna-Q $n{=}10$ is not above Q-learning). H2: on the blocking maze, Dyna-Q's post-change reward falls with $n$ and Dyna-Q+ recovers better. H3: on the shortcut maze, Dyna-Q+ gets more post-change reward than Dyna-Q. H4: under slip, the benefit of planning shrinks or reverses at large $n$. H5: prioritized sweeping matches Dyna-Q on static tasks and shares its staleness failure. The hypotheses and their refutation criteria were fixed before the final runs (BRIEF.md).

\textbf{Protocol.} Main comparison: $n=10$ for the three planning agents, 20 seeds (0--19), 5 tasks. Sweep over $n\in\{1,5,20,50,100\}$: 10 seeds (0--9), three planning agents, all tasks. Dyna-Q+ ablation: 20 seeds at $n=10$. Bonus sweep: $\kappa\in\{10^{-4},10^{-2},3{\cdot}10^{-2}\}$ plus the default $10^{-3}$, all on seeds 0--9. Hyperparameters ($\alpha=\rhval{const/alpha:1}$, $\gamma=\rhval{const/gamma:2}$, $\epsilon=\rhval{const/epsilon:1}$, $\kappa=10^{-3}$, $\theta=10^{-4}$) were fixed a priori and not tuned for any method, so there is no tuning budget and no held-out split; the environments are fixed and results are over seeds. The seed fixes both agent and environment randomness; Q-learning and prioritized sweeping share the random stream until the first reward, which explains identical early behaviour in some seeds. Tables show the mean over seeds, with the standard deviation where stated; figures show mean $\pm$ standard error. Where a sweep table lists Q-learning or Dyna-Q on seeds 0--9, these are the main-comparison runs of those seeds. Tests are two-sided Welch and paired-seed $t$-tests from the harness (\texttt{rh compare}), uncorrected for the many comparisons. Hardware: one CPU core of a shared Apple-silicon laptop-class machine, pure Python; each run takes under a second and the whole study a few minutes of wall-clock.
